\item En 1816 Robert Stirling, un clérigo escocés, patentó el \textit{motor Stirling}, que después encontró gran variedad de aplicaciones, incluyendo su empleo en colectores de energía solar para transformar la luz solar en electricidad. El combustible se quema externamente para calentar uno de los dos cilindros del motor. Una cantidad fija de gas inerte se mueve cíclicamente entre los cilindros y se expande en el caliente y se contrae en el frío. La figura P9.33 representa un modelo de su ciclo termodinámico. Considere $n$ moles de un gas ideal monoatómico que circula una vez por el ciclo, que consiste en dos procesos isotérmicos a temperaturas $3T_i$ y $T_i$ y dos procesos a volumen constante. Ahora el objetivo es obtener la eficiencia de esta máquina. 

(a) Encuentre la energía transferida por calor al gas durante el proceso isovolumétrico $AB$. 
(b) Determine la energía que se transfiere por calor al gas durante el proceso isotérmico $BC$. 
(c) Encuentre la energía transferida por calor al gas durante el proceso isovolumétrico $CD$. 
(d) Obtenga la energía que se transfiere por calor al gas durante el proceso isotérmico $DA$. 
(e) Identifique cuáles de los resultados de los incisos (a) a (d) son positivos y evalúe la energía de entrada a la máquina por calor. 
(f) De la primera ley de la termodinámica, encuentre el trabajo efectuado por la máquina. 
(g) De los resultados de los incisos (e) y (f), evalúe la eficiencia de la máquina. Un motor Stirling es más fácil de fabricar que un motor de combustión interna o una turbina. Puede funcionar con la quema de basura. Puede operar con la energía transferida por luz solar y no producir materiales de descarga. Los motores Stirling no se emplean en automóviles debido a su largo tiempo de arranque y a su mala respuesta de aceleración.
